% =============================================================================
% Chapter 2: Compute Infrastructure — ML Systems Lecture Slides (Volume II)
% =============================================================================
% This deck covers the physical infrastructure stack from accelerator to pod.
% Running example: 175B-parameter model on H100 cluster.
%
% IMAGES (SVG → PDF):
%   - accelerator-spectrum.pdf   - hbm-architecture.pdf
%   - roofline-h100.pdf          - bandwidth-hierarchy.pdf
%   - gpu-evolution.pdf          - infra-stack.pdf
%   - tco-overview.pdf
% =============================================================================
\documentclass[aspectratio=169, 12pt]{beamer}
\usepackage{../../assets/beamerthememlsys}

\mlsyssetup{
  volume       = {Volume II},
  chapter      = {Chapter 2},
  logo         = {../../assets/img/logo-mlsysbook.png},
  instlogo     = {../../assets/img/logo-harvard.png},
  chaptertitle = {Compute Infrastructure},
}

% --- Fonts ---
\usepackage{fontspec}
\setsansfont{Helvetica Neue}[
  BoldFont={Helvetica Neue Bold},
  ItalicFont={Helvetica Neue Italic},
  BoldItalicFont={Helvetica Neue Bold Italic},
]
\setmonofont{JetBrains Mono}[Scale=0.85]

% --- Packages ---
\usepackage{booktabs}
\usepackage{amsmath}

% --- Image paths ---
\graphicspath{
  {images/}
}

% --- Chapter-specific macros ---
\newcommand{\memwall}{\textcolor{computestroke}{\textbf{Memory Wall}}}
\newcommand{\powerwall}{\textcolor{errorstroke}{\textbf{Power Wall}}}
\newcommand{\commwall}{\textcolor{routingstroke}{\textbf{Communication Wall}}}

% --- Helper: safe image include ---
\newcommand{\safeimg}[2][width=\textwidth,keepaspectratio]{%
  \IfFileExists{images/#2}{\includegraphics[#1]{#2}}{%
    \IfFileExists{#2}{\includegraphics[#1]{#2}}{%
      \fbox{\parbox[c][2.5cm][c]{0.85\linewidth}{\centering\footnotesize\textcolor{midgray}{[Missing image]}}}%
    }%
  }%
}

% --- Section count for navigation ---
\setcounter{mlsystotalsections}{7}

\title{Compute Infrastructure}
\author{Vijay Janapa Reddi}
\institute{Harvard University}
\date{}

\begin{document}

% =============================================================================
% TITLE SLIDE
% =============================================================================
\mlsystitle{Compute Infrastructure}{The Physics of the ML Fleet}{cover_infrastructure.png}

% =============================================================================
% LEARNING OBJECTIVES
% =============================================================================
\begin{frame}{Learning Objectives}
\note{[2 min] Walk through objectives. Emphasize that this chapter maps the
\emph{physical} infrastructure from transistor to datacenter. Ask: ``How many
of you have thought about where the electricity comes from to train GPT-4?''}

\small
\begin{enumerate}
  \item Explain the \textbf{Accelerator Spectrum} (GPU vs.\ TPU vs.\ ASIC) and the \textbf{Generality Tax}
  \item Quantify the \textbf{Memory Wall} using HBM generations and token latency analysis
  \item Apply the \textbf{Roofline Model} to diagnose memory-bound vs.\ compute-bound workloads
  \item Analyze the \textbf{Bandwidth Hierarchy} (HBM $\to$ NVLink $\to$ InfiniBand) and its impact on parallelism
  \item Evaluate \textbf{Cooling Architectures} and calculate power requirements using \textbf{PUE}
  \item Perform a \textbf{TCO} analysis to determine break-even for build-vs-buy decisions
\end{enumerate}

\end{frame}

\begin{frame}{Visual Language}
\note{[1 min] Explain the semantic color system used throughout the course.
These colors are consistent across all diagrams and slides.}

\small
Throughout this course, colors carry meaning:

\vspace{0.3cm}
\begin{columns}[T]
  \begin{column}{0.45\textwidth}
    \begin{mlsyscard}{computestroke}
    \textbf{Blue} --- Compute / Processing\\
    {\footnotesize GPU ops, forward/backward pass, inference}
    \end{mlsyscard}
    \vspace{0.15cm}
    \begin{mlsyscard}{datastroke}
    \textbf{Green} --- Data / Memory\\
    {\footnotesize Data flow, caches, healthy paths}
    \end{mlsyscard}
  \end{column}
  \begin{column}{0.45\textwidth}
    \begin{mlsyscard}{routingstroke}
    \textbf{Orange} --- Routing / Scheduling\\
    {\footnotesize Load balancers, batch windows}
    \end{mlsyscard}
    \vspace{0.15cm}
    \begin{mlsyscard}{errorstroke}
    \textbf{Red} --- Error / Cost / Bottleneck\\
    {\footnotesize Loss, decode phase, waste}
    \end{mlsyscard}
  \end{column}
\end{columns}
\end{frame}



% =============================================================================
\section{Infra Stack}
% =============================================================================

\begin{frame}{The Infrastructure Stack}
\note{[3 min] Overview slide. The running example is a 175B model. Each level
of the stack exists because a constraint at the level below became intolerable.
Walk through bottom-to-top: accelerator (Memory Wall), node (capacity), rack
(Power Wall), pod (Communication Wall). Ask: ``Which wall do you think is
hardest to overcome?''}

% --- Layout: FULL-WIDTH diagram ---
\centering
\safeimg[width=0.92\textwidth,height=4.8cm,keepaspectratio]{infra-stack.pdf}

\vspace{0.1cm}
\mlsysinsight{Key principle:}{Each constraint at one level creates the engineering motivation for the next level.}

\end{frame}

\begin{frame}{Three Walls Across the Stack}
\note{[2 min] The three walls are not independent; they interact. Solving one
often shifts pressure to another. This is the C\textsuperscript{3} taxonomy
from Ch1 applied to physical infrastructure. If short: focus on the cascade
pattern, not individual walls.}

\small
\renewcommand{\arraystretch}{1.2}
{\footnotesize
\begin{tabular}{@{}llll@{}}
  \toprule
  \textbf{Level} & \textbf{\memwall} & \textbf{\powerwall} & \textbf{\commwall} \\
  \midrule
  \textbf{Accelerator} & HBM bandwidth & 700 W TDP & On-die mesh \\
  \textbf{Node}        & 640 GB aggregate & 7 kW per node & NVLink: 900 GB/s \\
  \textbf{Rack}        & PCIe to host DRAM & 33--40 kW/rack & ToR switch \\
  \textbf{Pod}         & Distributed storage & 7--10 MW & IB: 50 GB/s \\
  \bottomrule
\end{tabular}
}

\vspace{0.2cm}
\begin{mlsyscard}{crimson}
\textbf{Running Example}: 175B model $\to$ 350 GB weights (FP16) $\to$ 1.4 TB training state $\to$ does not fit on one chip.
\end{mlsyscard}

\end{frame}

% =============================================================================
\section{Accelerators}
% =============================================================================

\begin{frame}{The Accelerator Spectrum}
\note{[3 min] The spectrum is not a ranking; it is a trade-off continuum.
Key insight: the Generality Tax. A CPU spends 90\% of its die area on things
ML does not need. Each step right reclaims die area for arithmetic.
Ask: ``Why would anyone choose a GPU over a custom ASIC?''}

% --- Layout: FULL-WIDTH diagram ---
\centering
\safeimg[width=0.92\textwidth,height=4.8cm,keepaspectratio]{accelerator-spectrum.pdf}

\vspace{0.1cm}
\mlsysconcept{The Generality Tax:}{Every transistor spent on branch prediction is a transistor that could have been a multiplier.}

\end{frame}

\begin{frame}{GPU Architecture Evolution}
\note{[3 min] Walk through four generations. The pattern: each generation
addresses the bottleneck exposed by the previous one. Volta: compute. Ampere:
precision. Hopper: Transformers. Blackwell: die size limit.
Key number: efficiency grew 70$\times$ but TFLOPS grew 236$\times$ ---
the power grid must absorb the gap.}

% --- Layout: FULL-WIDTH timeline ---
\centering
\safeimg[width=0.92\textwidth,height=4.8cm,keepaspectratio]{gpu-evolution.pdf}

\vspace{0.1cm}
\mlsysalert{Pattern:}{Each generation solves the previous generation's bottleneck. Hardware lifetime is shrinking.}

\end{frame}

\begin{frame}{GPU Specs: Four Generations}
\note{[2 min] Reference table. Emphasize that NVLink BW doubles each generation
--- tracking model size growth. The TFLOPS/W column reveals the efficiency wall.
If short: just highlight the 3 key columns.}

\small
\renewcommand{\arraystretch}{1.15}
{\footnotesize
\begin{tabular}{@{}llrrrr@{}}
  \toprule
  \textbf{Arch.} & \textbf{Year} & \textbf{Peak TF} & \textbf{TDP} & \textbf{NVLink} & \textbf{TF/W} \\
  \midrule
  Volta (V100)     & 2017 &  125 & 300 W &  300 GB/s & 0.42 \\
  Ampere (A100)    & 2020 &  312 & 400 W &  600 GB/s & 0.78 \\
  Hopper (H100)    & 2022 & 1,979 & 700 W &  900 GB/s & 2.83 \\
  Blackwell (B200) & 2024 & 4,500 & 1,000 W & 1,800 GB/s & 4.50 \\
  \bottomrule
\end{tabular}
}

\vspace{0.15cm}
\begin{mlsyscard}{errorstroke}
{\footnotesize \textbf{The Efficiency Wall}: Raw TFLOPS grew \textbf{236$\times$} (P100$\to$B200), but TFLOPS/Watt grew only \textbf{70$\times$}. The gap is what cooling and power infrastructure must absorb.}
\end{mlsyscard}

\end{frame}

% =============================================================================
\section{Memory Wall}
% =============================================================================

\begin{frame}{The Memory Wall}
\note{[3 min] This is the central constraint for inference. Even with 2,000
TFLOPS, the H100's Tensor Cores sit idle 99\% of the time during single-token
generation. The bottleneck is delivery, not computation.
Ask: ``If you could double either FLOPS or bandwidth, which helps more for
serving a chatbot?''}

\small
\begin{columns}[T]
  \begin{column}{0.52\textwidth}
    \safeimg[width=\textwidth,height=5.5cm,keepaspectratio]{hbm-architecture.pdf}
  \end{column}
  \begin{column}{0.45\textwidth}
    \textbf{The paradox:} faster logic idles more, waiting for data.

    \vspace{0.15cm}
    {\footnotesize
    \renewcommand{\arraystretch}{1.15}
    \begin{tabular}{@{}lr@{}}
      \toprule
      \textbf{Metric} & \textbf{Value} \\
      \midrule
      DDR5 BW & $\sim$200 GB/s \\
      HBM3 BW & \textbf{3,350 GB/s} \\
      \textbf{Ratio} & \textbf{50$\times$} \\
      \midrule
      DDR5 energy & 20 pJ/bit \\
      HBM3 energy & \textbf{2 pJ/bit} \\
      \bottomrule
    \end{tabular}
    }

    \vspace{0.1cm}
    \mlsysalert{Key:}{HBM \emph{moves} the Memory Wall but does not eliminate it.}
  \end{column}
\end{columns}

\end{frame}

\begin{frame}{Token Latency: The Physics}
\note{[3 min] Worked example. Llama-3 70B on H100. Walk through the two
calculations step by step. The reveal: 97\% of time is memory transfer.
Even infinite compute would not help. This is why HBM bandwidth improvements
deliver nearly linear inference speedups.}

\small
\textbf{Scenario: One token of Llama-3 70B on H100}

\vspace{0.1cm}
\renewcommand{\arraystretch}{1.2}
{\footnotesize
\begin{tabular}{@{}llll@{}}
  \toprule
  \textbf{Component} & \textbf{Calculation} & \textbf{Time} & \textbf{Share} \\
  \midrule
  \textcolor{computestroke}{Compute} & $2 \times 70\text{B} / 1{,}979\text{ TF}$ & \textbf{0.07 ms} & 3\% \\
  \textcolor{datastroke}{Memory}  & $140\text{ GB} / 3.35\text{ TB/s}$ & \textbf{41.8 ms} & \textbf{97\%} \\
  \midrule
  & & $T_{\text{token}} \approx$ & \textbf{41.9 ms} \\
  \bottomrule
\end{tabular}
}

\vspace{0.15cm}
\begin{mlsyscard}{errorstroke}
\textbf{The processor spends 97\% of its time waiting for data.}
Even a hypothetical chip with \emph{infinite} compute would generate tokens only negligibly faster. Memory bandwidth is the binding constraint.
\end{mlsyscard}

\end{frame}

% --- ACTIVE LEARNING 1: Predict ---
\begin{frame}{Predict: Memory-Bound or Compute-Bound?}
\note{[2 min] Prediction exercise before revealing the Roofline Model.
Give students 60 seconds. Do NOT reveal the framework yet.
Ask 2--3 students to share their reasoning.}

\centering
\vspace{0.8cm}
{\Large\bfseries Think--Write--Share}

\vspace{0.5cm}
{\large A team trains a 175B model at batch size 2048.\\
Another team serves the same model at batch size 1.\\[0.3cm]
\alert{Which workload benefits more from buying\\a faster GPU (more TFLOPS)?}}

\vspace{0.5cm}
{\normalsize Write your answer and reasoning. \textcolor{midgray}{(60 seconds)}}

\end{frame}

% =============================================================================
\section{Roofline}
% =============================================================================

\mlsysfocus{The Roofline Model}{%
$\text{Achievable FLOPS} = \min\left(R_{\text{peak}},\; \text{BW} \times I\right)$\\[0.3cm]
\normalsize where $I = \text{FLOP/byte}$ (arithmetic intensity)%
}

\begin{frame}{Roofline: H100 Landscape}
\note{[3 min] Reveal after the prediction exercise. Walk through the plot.
The ridge point at $\sim$591 FLOP/byte separates the two regimes.
LLM decode at B=1 achieves less than 1\% of peak. LLM training crosses the
ridge point. The answer to the prediction: training benefits from more TFLOPS;
inference benefits from more bandwidth.}

% --- Layout: FULL-WIDTH diagram ---
\centering
\safeimg[width=0.92\textwidth,height=4.8cm,keepaspectratio]{roofline-h100.pdf}

\vspace{0.1cm}
\mlsysconcept{Diagnostic power:}{The Roofline tells you \emph{which resource to optimize} before you spend money.}

\end{frame}

\begin{frame}{Roofline: Workload Placement}
\note{[2 min] Expand on how batch size shifts arithmetic intensity. Doubling
batch size doubles FLOP for same bytes transferred $\to$ shifts workload right.
Key insight: same hardware, same model, but different operating regimes
depending on batch size. This is why training and inference have different
optimal hardware.}

\footnotesize
\renewcommand{\arraystretch}{1.15}
\begin{tabular}{@{}lrrl@{}}
  \toprule
  \textbf{Workload} & \textbf{$I$ (FLOP/B)} & \textbf{\% Peak} & \textbf{Regime} \\
  \midrule
  LLM decode (B=1)    & $\sim$1    & $<$1\%  & \textcolor{computestroke}{\textbf{Memory-bound}} \\
  LLM decode (B=32)   & $\sim$32   & $\sim$5\%  & \textcolor{computestroke}{Memory-bound} \\
  CNN training         & 50--200    & 30--60\% & Near ridge \\
  LLM prefill          & 100--500   & 50--80\% & Transitioning \\
  LLM training (175B)  & $\sim$2000 & 30--50\% & \textcolor{routingstroke}{\textbf{Compute-bound}} \\
  \bottomrule
\end{tabular}

\vspace{0.15cm}
\begin{columns}[T]
  \begin{column}{0.48\textwidth}
    \begin{mlsyscard}{computestroke}
    {\footnotesize \textbf{Below ridge:} Buy bandwidth.\\
    Faster GPU is wasted money.}
    \end{mlsyscard}
  \end{column}
  \begin{column}{0.48\textwidth}
    \begin{mlsyscard}{routingstroke}
    {\footnotesize \textbf{Above ridge:} Buy compute.\\
    Faster HBM is wasted money.}
    \end{mlsyscard}
  \end{column}
\end{columns}

\end{frame}

% --- ACTIVE LEARNING 2: Exercise ---
\begin{frame}{Your Turn: Roofline Diagnosis}
\note{[3 min] Give students 90 seconds. Answer: $I = 2 \times 7\text{B} \times 2 / (14\text{ GB}) = 2$ FLOP/byte.
Ridge point $\approx$ 591. Since 2 $\ll$ 591, deeply memory-bound.
Adding more TFLOPS yields zero benefit. Need more bandwidth.}

\small
\begin{columns}[T]
  \begin{column}{0.55\textwidth}
    {\normalsize\bfseries Diagnose the Bottleneck}

    \vspace{0.2cm}
    A 7B model (14 GB FP16 weights) served at batch size 1 on an H100:
    \begin{itemize}
      \item $R_{\text{peak}}$ = 1,979 TFLOPS
      \item BW = 3.35 TB/s
      \item Ridge point $\approx$ 591 FLOP/byte
    \end{itemize}

    \vspace{0.1cm}
    \textbf{Calculate} the arithmetic intensity and determine the regime.

    {\footnotesize\textcolor{midgray}{(90 seconds --- then compare with a neighbor)}}
  \end{column}
  \begin{column}{0.42\textwidth}
    \pause
    \begin{mlsyscard}{datastroke}
    \textbf{Solution:}\\[0.1cm]
    {\footnotesize
    $I = \frac{2 \times 7\text{B}}{14\text{ GB}} = 1$ FLOP/byte\\[0.1cm]
    Ridge point = 591\\
    $1 \ll 591$\\[0.1cm]
    \textcolor{computestroke}{\textbf{Deeply memory-bound!}}\\[0.05cm]
    More TFLOPS $\to$ zero benefit.\\
    More BW $\to$ linear speedup.
    }
    \end{mlsyscard}
  \end{column}
\end{columns}

\end{frame}

% =============================================================================
\section{Node \& BW}
% =============================================================================

\begin{frame}{The Node: Aggregating Accelerators}
\note{[3 min] A 175B model needs 350 GB for weights alone. One H100 has 80 GB.
We need at least 5 GPUs for weights, and with optimizer state (1.4 TB total),
all 8 in a DGX node are barely sufficient. The node bridges the capacity gap
by connecting 8 GPUs with NVLink. Key: 8$\times$ capacity, 8$\times$ compute.}

\small
\begin{columns}[T]
  \begin{column}{0.52\textwidth}
    \textbf{Why nodes exist:} HBM has bandwidth but not enough \emph{capacity}.

    \vspace{0.15cm}
    {\footnotesize
    \renewcommand{\arraystretch}{1.15}
    \begin{tabular}{@{}lr@{}}
      \toprule
      \textbf{Component} & \textbf{175B Model} \\
      \midrule
      Weights (FP16)     & 350 GB \\
      Gradients          & 350 GB \\
      Optimizer (Adam)   & 700 GB \\
      \textbf{Total}     & \textbf{1.4 TB} \\
      \midrule
      Single H100 HBM    & 80 GB \\
      8-GPU Node HBM     & \textbf{640 GB} \\
      \bottomrule
    \end{tabular}
    }

    \vspace{0.1cm}
    {\footnotesize Still not enough for full training state $\to$ ZeRO sharding, offloading.}
  \end{column}
  \begin{column}{0.45\textwidth}
    \begin{mlsyscard}{crimson}
    \textbf{DGX H100 Node}\\[0.1cm]
    {\footnotesize
    8$\times$ H100 SXM GPUs\\
    640 GB aggregate HBM\\
    NVSwitch: 900 GB/s per GPU\\
    4$\times$ NVSwitch chips\\
    Total: $\sim$7 kW}
    \end{mlsyscard}

    \vspace{0.1cm}
    \begin{mlsyscard}{datastroke}
    {\footnotesize \textbf{Key insight:} NVSwitch transforms 8 independent GPUs into one ``virtual accelerator'' with unified memory.}
    \end{mlsyscard}
  \end{column}
\end{columns}

\end{frame}

\begin{frame}{The Bandwidth Hierarchy}
\note{[3 min] This is the most important systems diagram in the chapter.
Each zone boundary drops bandwidth by $\sim$10$\times$. The 18$\times$ cliff
between NVLink and InfiniBand dictates that tensor parallelism \emph{must}
stay within a node. Data parallelism can span the fabric because it
communicates less frequently.}

% --- Layout: FULL-WIDTH pyramid ---
\centering
\safeimg[width=0.88\textwidth,height=4.5cm,keepaspectratio]{bandwidth-hierarchy.pdf}

\vspace{0.1cm}
\small
\begin{columns}[T]
  \begin{column}{0.30\textwidth}
    \centering
    \textcolor{computestroke}{\textbf{Tensor Parallel}}\\
    {\scriptsize Within node (NVLink)}
  \end{column}
  \begin{column}{0.30\textwidth}
    \centering
    \textcolor{routingstroke}{\textbf{Pipeline Parallel}}\\
    {\scriptsize Across nodes (IB)}
  \end{column}
  \begin{column}{0.30\textwidth}
    \centering
    \textcolor{datastroke}{\textbf{Data Parallel}}\\
    {\scriptsize Across pods (fabric)}
  \end{column}
\end{columns}

\end{frame}

\begin{frame}{Bandwidth Hierarchy: The Numbers}
\note{[2 min] Reference table. The 18$\times$ cliff is the critical number.
Emphasize that these are not engineering failures to fix; they reflect
physics of each interconnect medium. If short: just highlight the NVLink
to IB transition.}

\footnotesize
\renewcommand{\arraystretch}{1.15}
\begin{tabular}{@{}llrrl@{}}
  \toprule
  \textbf{Domain} & \textbf{Interconnect} & \textbf{BW} & \textbf{Latency} & \textbf{Scope} \\
  \midrule
  Intra-Package   & Si Interposer  & $\sim$3,350 GB/s & $<$100 ns  & Single chip \\
  Intra-Node      & NVLink 4.0     & $\sim$900 GB/s   & $\sim$1 $\mu$s  & 8 GPUs \\
  Intra-Node (IO) & PCIe Gen5 x16  & $\sim$64 GB/s    & $\sim$2 $\mu$s  & CPU--GPU \\
  Inter-Node      & IB NDR 400G    & $\sim$50 GB/s    & 5--10 $\mu$s & 1000s GPUs \\
  \bottomrule
\end{tabular}

\vspace{0.15cm}
\begin{mlsyscard}{errorstroke}
{\footnotesize \textbf{The 18$\times$ cliff:} NVLink (900 GB/s) $\to$ InfiniBand (50 GB/s). This single boundary dictates the entire parallelism hierarchy for distributed training.}
\end{mlsyscard}

\end{frame}

% --- ACTIVE LEARNING: Micro-Retrieval Cue ---
\begin{frame}{Quick Check}
\note{[1 min] Answer: NVLink 900 GB/s vs IB 50 GB/s. TP must stay intra-node.}

\centering
\vspace{1.0cm}
{\Large\bfseries Quick Check}

\vspace{0.6cm}
{\large What is the 18$\\times$ bandwidth cliff,\\and why does it matter for parallelism?}

\vspace{0.5cm}
{\normalsize\textcolor{midgray}{15 seconds --- then I will cold-call.}}

\end{frame}



% =============================================================================
\section{Rack \& Pod}
% =============================================================================

\begin{frame}{The Rack: Power and Cooling}
\note{[3 min] A traditional rack draws 5--10 kW. An ML rack with 4 DGX H100
nodes draws 33--40 kW. This is an industrial furnace, not a server closet.
Two key concepts: power delivery chain (5 stages, 85--90\% cumulative
efficiency) and PUE. Ask: ``Why is liquid cooling a necessity, not luxury?''}

\small
\begin{columns}[T]
  \begin{column}{0.52\textwidth}
    \textbf{ML rack power density:} 10$\times$ traditional DC

    \vspace{0.15cm}
    {\footnotesize
    \renewcommand{\arraystretch}{1.15}
    \begin{tabular}{@{}lr@{}}
      \toprule
      \textbf{Component} & \textbf{Power} \\
      \midrule
      32 GPUs ($\times$700 W)  & 22.4 kW \\
      CPUs, NVSwitch, NICs      & $\sim$8 kW \\
      Power conversion losses   & $\sim$3 kW \\
      \textbf{Total per rack}   & \textbf{33--40 kW} \\
      \midrule
      Traditional server rack   & 5--10 kW \\
      \bottomrule
    \end{tabular}
    }

    \vspace{0.1cm}
    {\footnotesize 128 racks for 175B training $\to$ $\sim$5 MW cluster.}
  \end{column}
  \begin{column}{0.45\textwidth}
    \begin{mlsyscard}{routingstroke}
    \textbf{PUE: Cooling Tax}\\[0.1cm]
    {\footnotesize
    $\text{PUE} = \frac{\text{Total Facility Power}}{\text{IT Equipment Power}}$\\[0.1cm]
    Air-cooled: PUE 1.4--1.6\\
    (40--60\% waste!)\\[0.05cm]
    \textbf{Liquid-cooled: PUE 1.05--1.10}\\
    (5--10\% overhead only)}
    \end{mlsyscard}

    \vspace{0.1cm}
    \begin{mlsyscard}{errorstroke}
    {\footnotesize For 1,000 GPUs: liquid cooling saves \textbf{294 kW} $\to$ \$180K/yr in electricity.}
    \end{mlsyscard}
  \end{column}
\end{columns}

\end{frame}

\begin{frame}{The Pod: Warehouse-Scale Computer}
\note{[3 min] The pod is where the Communication Wall hits at full force.
1,024 DGX H100 nodes = 8,192 GPUs, 250 racks, 7--10 MW. Key math:
physics-limit training time for 175B is $\sim$12 hours. Actual time is 2--4
weeks due to overhead. The 15--25$\times$ gap is the central challenge of
distributed systems engineering.}

\small
\begin{columns}[T]
  \begin{column}{0.55\textwidth}
    \textbf{1,024-node DGX H100 cluster:}

    \vspace{0.1cm}
    {\footnotesize
    \renewcommand{\arraystretch}{1.1}
    \begin{tabular}{@{}lr@{}}
      \toprule
      \textbf{Resource} & \textbf{Scale} \\
      \midrule
      GPUs            & 8,192 \\
      Racks           & $\sim$250 \\
      Power draw      & 7--10 MW \\
      Weight          & 300+ metric tons \\
      IB switches     & Hundreds \\
      Cables          & Tens of thousands \\
      \bottomrule
    \end{tabular}
    }

    \vspace{0.1cm}
    {\footnotesize The datacenter floor must support 1,500--2,500 kg/m$^2$ (vs.\ 500--1,000 traditional).}
  \end{column}
  \begin{column}{0.42\textwidth}
    \begin{mlsyscard}{crimson}
    \textbf{175B Training Time}\\[0.1cm]
    {\footnotesize
    Physics limit: $\sim$\textbf{12 hours}\\
    {\scriptsize (8,192 GPUs at 45\% MFU)}\\[0.05cm]
    Actual: \textbf{2--4 weeks}\\
    {\scriptsize (comm, bubbles, checkpoints, failures)}\\[0.1cm]
    Gap: \textbf{15--25$\times$}\\
    Closing this gap is distributed\\
    systems engineering.}
    \end{mlsyscard}
  \end{column}
\end{columns}

\end{frame}

\begin{frame}{Scaling Efficiency}
\note{[2 min] As we add more GPUs, communication overhead grows. Scaling
efficiency is the fraction of ideal speedup actually achieved. The formula
quantifies the diminishing returns. Key: beyond 2,000--4,000 GPUs, marginal
benefit per GPU approaches zero for most model sizes.}

\small
$$\eta_{\text{scale}} = \frac{T_{\text{compute}}}{T_{\text{compute}} + T_{\text{comm}} + T_{\text{idle}}}$$

\vspace{0.1cm}
\renewcommand{\arraystretch}{1.15}
{\footnotesize
\begin{tabular}{@{}rrrl@{}}
  \toprule
  \textbf{GPUs} & \textbf{Ideal Speedup} & \textbf{Actual} & \textbf{$\eta_{\text{scale}}$} \\
  \midrule
  8     &  8$\times$    & 7.6$\times$   & 95\% \\
  64    &  64$\times$   & 54$\times$    & 85\% \\
  1,024 & 1,024$\times$ & 615$\times$   & 60\% \\
  8,192 & 8,192$\times$ & 3,277$\times$ & 40\% \\
  \bottomrule
\end{tabular}
}

\vspace{0.1cm}
\mlsysalert{Amdahl at scale:}{The first 100 GPUs provide near-linear speedup. Beyond 2,000--4,000, marginal benefit per GPU approaches zero.}

\end{frame}

% --- ACTIVE LEARNING 3: Discussion ---
\begin{frame}{Discussion: Which Infrastructure Fails First?}
\note{[3 min] Turn-and-talk. No single right answer. The point: infrastructure
constraints are interdependent. A team might over-provision GPUs and under-provision
cooling. Cold-call 2--3 pairs.}

\centering
\vspace{0.6cm}
{\large\bfseries Turn and Talk \textcolor{midgray}{(90 seconds)}}

\vspace{0.5cm}
{\large A startup orders 256 H100 GPUs to train\\a 70B model. They arrive in 6 months.\\[0.3cm]
\alert{What infrastructure constraint is most likely\\to delay their first training run?}}

\vspace{0.5cm}
\begin{columns}[c]
  \begin{column}{0.18\textwidth}\centering\small Power\\ Delivery\end{column}
  \begin{column}{0.18\textwidth}\centering\small Cooling\\ Capacity\end{column}
  \begin{column}{0.18\textwidth}\centering\small Network\\ Fabric\end{column}
  \begin{column}{0.18\textwidth}\centering\small Facility\\ Space\end{column}
  \begin{column}{0.18\textwidth}\centering\small Last Mile\\ (Burn-in)\end{column}
\end{columns}

\end{frame}

% =============================================================================
\section{TCO}
% =============================================================================

\begin{frame}{Total Cost of Ownership}
\note{[3 min] TCO is the framework that connects all physical constraints to
dollars. The key insight: utilization is the single variable that flips the
build-vs-buy decision. Same cluster, same facility --- brilliant at 80\% util,
financial disaster at 20\% util.}

% --- Layout: FULL-WIDTH diagram ---
\centering
\safeimg[width=0.92\textwidth,height=4.8cm,keepaspectratio]{tco-overview.pdf}

\vspace{0.1cm}
\mlsysconcept{Central tension:}{Utilization is the single variable that flips the build-vs-buy decision.}

\end{frame}

\begin{frame}{TCO: The Numbers}
\note{[2 min] Walk through the three key numbers. On-prem at 80\% util is
$\sim$\$2.40/GPU-hour. Cloud on-demand is \$4.00. Break-even at $\sim$55\%
utilization. The comparison changes with reserved instances (40--60\% discount).
If short: focus on the break-even number.}

\small
\renewcommand{\arraystretch}{1.15}
{\footnotesize
\begin{tabular}{@{}lll@{}}
  \toprule
  & \textbf{On-Premises} & \textbf{Cloud (On-Demand)} \\
  \midrule
  \textbf{CapEx}       & \$350K/node (amortize 3 yr) & \$0 \\
  \textbf{Electricity} & $\sim$\$5K/yr/node          & Included \\
  \textbf{Annual cost} & $\sim$\$122K/node (80\% util) & $\sim$\$224K/node (80\% util) \\
  \textbf{Eff.\ rate}  & $\sim$\$2.40/GPU-hour       & \$4.00/GPU-hour \\
  \midrule
  \textbf{Break-even}  & \multicolumn{2}{c}{\alert{$\sim$55\% sustained utilization}} \\
  \bottomrule
\end{tabular}
}

\vspace{0.15cm}
\begin{columns}[T]
  \begin{column}{0.48\textwidth}
    \begin{mlsyscard}{datastroke}
    {\footnotesize \textbf{Own if}: continuous 24/7 workloads, predictable demand, 3+ year horizon.}
    \end{mlsyscard}
  \end{column}
  \begin{column}{0.48\textwidth}
    \begin{mlsyscard}{errorstroke}
    {\footnotesize \textbf{Rent if}: bursty demand, rapid experimentation, $<$2 year horizon.}
    \end{mlsyscard}
  \end{column}
\end{columns}

\end{frame}

% --- ACTIVE LEARNING 4 (Exercise): Build vs Buy ---
\begin{frame}{Your Turn: Build vs.\ Buy}
\note{[3 min] Give students 90 seconds. Answer: Cloud cost = 500 GPUs $\times$
\$4 $\times$ 24 $\times$ 21 = \$1,008,000 per run. 4 runs = \$4.03M per year.
On-prem: 63 nodes $\times$ \$350K / 3 + electricity = \$7.8M + \$315K =
$\sim$\$8.1M/yr. Cloud is cheaper at 4 runs/yr. Need $\sim$10+ runs/yr for
on-prem to win.}

\small
\begin{columns}[T]
  \begin{column}{0.55\textwidth}
    {\normalsize\bfseries Build vs.\ Buy Decision}

    \vspace{0.15cm}
    Your team plans to train a 70B model:
    \begin{itemize}
      \item 500 H100 GPUs for 3 weeks per run
      \item 4 training runs per year
      \item Cloud rate: \$4.00/GPU-hour
      \item On-prem: \$350K/node, 3-yr amortization
    \end{itemize}

    \vspace{0.1cm}
    \textbf{Calculate} annual cost for both options.\\
    \textbf{How many runs/year} to justify on-prem?

    {\footnotesize\textcolor{midgray}{(90 seconds)}}
  \end{column}
  \begin{column}{0.42\textwidth}
    \pause
    \begin{mlsyscard}{datastroke}
    \textbf{Solution:}\\[0.05cm]
    {\footnotesize
    Cloud: 500$\times$\$4$\times$504h = \$1.01M/run\\
    4 runs = \textbf{\$4.03M/yr}\\[0.1cm]
    On-prem: 63 nodes $\times$ \$117K/yr\\
    + elec = \textbf{$\sim$\$7.7M/yr}\\[0.1cm]
    Break-even: $\sim$8 runs/yr\\[0.05cm]
    \textcolor{errorstroke}{\textbf{At 4 runs: cloud wins.}}\\
    \textcolor{datastroke}{\textbf{At 10+ runs: on-prem wins.}}
    }
    \end{mlsyscard}
  \end{column}
\end{columns}

\end{frame}

% =============================================================================
\section{Wrap-Up}
% =============================================================================

\begin{frame}{Fallacies}
\note{[2 min] Three common misconceptions with quantitative evidence.}

\small
\textbf{Fallacy:} \textit{More GPUs always means faster training.}\\
{\footnotesize Communication overhead grows with cluster size. Scaling from 1,024 to 2,048 GPUs doubles cost but may reduce training time by only 30--40\%.}

\vspace{0.15cm}
\textbf{Fallacy:} \textit{Peak FLOPS determines training throughput.}\\
{\footnotesize LLM inference operates at 1--2 FLOP/byte, achieving $<$1\% of peak FLOPS. The Roofline Model's ridge point is the diagnostic that matters, not the spec sheet.}

\vspace{0.15cm}
\textbf{Fallacy:} \textit{Liquid cooling is too expensive.}\\
{\footnotesize Over 3 years, a 1,000-GPU cluster saves \$540K in electricity with liquid cooling --- typically exceeding the incremental CapEx.}

\end{frame}

\begin{frame}{Pitfalls}
\note{[2 min] Three operational pitfalls. The power/cooling pitfall is the
most common: teams buy GPUs before confirming the facility can handle them.}

\footnotesize
\textbf{Pitfall:} \textit{Ignoring power and cooling constraints during planning.}\\
A rack of 4 DGX H100s requires 33--40 kW. Upgrading facility electrical infrastructure takes 12--18 months. GPUs arrive in 6 months. Result: expensive silicon in shipping containers in the parking lot.

\vspace{0.15cm}
\textbf{Pitfall:} \textit{Buying the latest accelerator without checking arithmetic intensity.}\\
B200 (4,500 TF) for 7B inference at B=1: arithmetic intensity $\approx$ 1 FLOP/byte. The B200 runs only 2.4$\times$ faster than H100 despite 2.3$\times$ higher peak FLOPS. You paid for compute you cannot use.

\vspace{0.15cm}
\textbf{Pitfall:} \textit{Treating all GPU-hours as equivalent in TCO analysis.}\\
An H100 GPU-hour delivers 6.3$\times$ the TFLOPS of an A100 hour. The right metric is cost per useful \textbf{TFLOP-hour}, not cost per GPU-hour.

\end{frame}

% --- RETRIEVAL PRACTICE ---

% --- MUDDIEST POINT ---
\begin{frame}{Muddiest Point}
\note{[2 min] Quick anonymous poll. Students write on a slip of paper or submit
digitally. Collect and scan for patterns. Address the top 2--3 confusions in the
next lecture's opening. This closes the feedback loop.}

\centering
\vspace{1.0cm}
{\Large\bfseries What was the \alert{muddiest point} today?}

\vspace{0.8cm}
{\normalsize Write down the concept you found \textbf{most confusing}.}

\vspace{0.5cm}
{\small\textcolor{midgray}{Anonymous. One sentence. Submit before you leave.}}

\end{frame}

\begin{frame}{What Were the Key Ideas?}
\note{[2 min] Retrieval practice. Students write 90 seconds, no notes.
Walk around the room. Do NOT show next slide yet.}

\centering
\vspace{1.5cm}
{\Large\bfseries Close your notes.}

\vspace{0.8cm}
{\large Write down the \textbf{4 most important concepts} from today.}

\vspace{0.8cm}
{\normalsize\textcolor{midgray}{90 seconds --- no peeking.}}

\end{frame}

\begin{frame}{Key Takeaways}
\note{[2 min] Reveal. Walk through each bullet. Emphasize the quantitative
anchors: 97\% memory time, 18$\times$ bandwidth cliff, PUE 1.05, 55\%
break-even utilization. These numbers calibrate engineering intuition.}

\scriptsize
\begin{itemize}\setlength\itemsep{0pt}
  \item \textbf{Generality Tax}: CPUs devote 5--10\% of die to math; GPUs 50--60\%; ASICs 90\%+. Flexibility traded for throughput.
  \item \textbf{Memory Wall}: Token generation: 97\% memory transfer, 3\% compute. HBM: 50$\times$ DDR BW, 10$\times$ efficiency.
  \item \textbf{Roofline Model}: Ridge point = $R_{\text{peak}}/\text{BW}$. Below $\to$ bandwidth-bound. Above $\to$ compute-bound. Diagnose before optimizing.
  \item \textbf{Bandwidth Hierarchy}: 18$\times$ cliff (NVLink $\to$ IB) dictates parallelism. Tensor parallel within node; data parallel across.
  \item \textbf{Power \& Cooling}: ML racks: 33--40 kW (10$\times$ traditional). Liquid PUE 1.05 vs.\ air PUE 1.4--1.6.
  \item \textbf{TCO}: Build-vs-buy break-even at $\sim$55\% utilization. Metric: cost/TFLOP-hour, not cost/GPU-hour.
  \item \textbf{Stack}: Each level's constraint creates the next level. No component optimized in isolation.
\end{itemize}

\end{frame}

\begin{frame}{References}
\note{[1 min] Point students to canonical papers.}

\small
\mlsysref{Jouppi+17}{Jouppi et al. ``In-Datacenter Performance Analysis of a Tensor Processing Unit.'' ISCA 2017.}
\mlsysref{Williams+09}{Williams, Waterman, Patterson. ``Roofline: An Insightful Visual Performance Model.'' CACM 2009.}
\mlsysref{Barroso+18}{Barroso, Clidaras, H\"{o}lzle. \textit{The Datacenter as a Computer}. 3rd ed.\ Morgan \& Claypool 2018.}
\mlsysref{Patterson+21}{Patterson et al. ``Carbon Emissions and Large Neural Networks.'' 2021.}
\mlsysref{Narayanan+21}{Narayanan et al. ``Efficient Large-Scale Language Model Training on GPU Clusters.'' SC 2021.}

\end{frame}

\begin{frame}{Next Lecture: Network Fabrics}
\note{[1 min] Forward hook. The network is not a passive connector; it is an
active determinant of system performance. The 18$\times$ bandwidth cliff
is the central constraint that Ch3 addresses. RDMA, collective communication,
and topology all aim to flatten this hierarchy.}

\footnotesize
\begin{columns}[T]
  \begin{column}{0.30\textwidth}
    \centering
    \begin{mlsyscard}{computestroke}
    {\large\bfseries RDMA}\\[0.1cm]
    {\footnotesize Kernel-bypass\\data transfer}
    \end{mlsyscard}
  \end{column}
  \begin{column}{0.30\textwidth}
    \centering
    \begin{mlsyscard}{routingstroke}
    {\large\bfseries Collectives}\\[0.1cm]
    {\footnotesize AllReduce,\\All-to-All}
    \end{mlsyscard}
  \end{column}
  \begin{column}{0.30\textwidth}
    \centering
    \begin{mlsyscard}{datastroke}
    {\large\bfseries Topology}\\[0.1cm]
    {\footnotesize Fat-tree, torus,\\rail-optimized}
    \end{mlsyscard}
  \end{column}
\end{columns}

\vspace{0.3cm}
\centering
{\normalsize The network transforms a warehouse of expensive space heaters\\into a unified training system.}\\[0.1cm]
{\small\textbf{The 18$\times$ bandwidth cliff is the constraint Ch3 addresses.}}

\end{frame}



\appendix

\begin{frame}{Backup: Extended Roofline Examples}
\note{Backup: More roofline examples for different workloads.}

\footnotesize
\renewcommand{\arraystretch}{1.15}
\begin{tabular}{@{}lrrrl@{}}
  \toprule
  \textbf{Workload} & \textbf{AI} & \textbf{H100 \%} & \textbf{B200 \%} & \textbf{Regime Shift?} \\
  \midrule
  Embedding lookup  & 0.5   & $<$0.1\% & $<$0.1\% & No (always mem-bound) \\
  Attention (B=1)   & 2     & 0.7\%    & 0.4\%    & No \\
  Conv2D (ResNet)   & 150   & 51\%     & 27\%     & \textcolor{errorstroke}{\textbf{Yes}} \\
  GEMM (4096$^2$)   & 1,365 & 100\%    & 100\%    & No (always compute-bound) \\
  \bottomrule
\end{tabular}

\vspace{0.15cm}
{\footnotesize Conv2D shifts from compute-bound on H100 to memory-bound on B200 as the ridge point rises.}

\end{frame}

\begin{frame}{Backup: HBM Generations Compared}
\note{Backup: Detailed HBM comparison across generations.}

\footnotesize
\renewcommand{\arraystretch}{1.15}
\begin{tabular}{@{}lrrrr@{}}
  \toprule
  \textbf{Generation} & \textbf{Capacity} & \textbf{BW} & \textbf{Stacks} & \textbf{pJ/bit} \\
  \midrule
  HBM2   & 16 GB  & 900 GB/s   & 4 & 3.9 \\
  HBM2e  & 32 GB  & 1,200 GB/s & 4 & 3.0 \\
  HBM3   & 48 GB  & 2,400 GB/s & 6 & 2.5 \\
  HBM3e  & 80 GB  & 3,350 GB/s & 8 & 2.0 \\
  \bottomrule
\end{tabular}

\vspace{0.15cm}
{\footnotesize Each generation roughly doubles bandwidth while improving energy efficiency by 20--30\%.}

\end{frame}


\end{document}
